
% --- BEGIN USER-REQUESTED LITERATURE REVIEW STRUCTURE ---
\section{Literature Survey}

\subsection{Introduction}
The integration of robotics and artificial intelligence into healthcare is advancing rapidly, driven by challenges such as aging populations, rising healthcare costs, and medical staff shortages. Researchers have developed various robotic solutions, including telepresence robots for remote consultations, autonomous robots for logistics, and socially assistive robots (SARs) for patient engagement. While each type has proven useful, they often function in isolation, addressing only part of patients’ and clinicians’ overall needs. The goal is to show that although the components for such a system already exist, combining them into a single, user-friendly, and autonomous platform remains an underexplored area with high potential, especially for improving healthcare accessibility in regions like Oman.

Telepresence robots, often described as a mobile screen, are widely used for remote medical consultations. Systems such as InTouch Health RP-VITA and the VGo robot have been used successfully in “telestroke” care, allowing neurologists to assess patients remotely, reduce diagnosis time, and improve outcomes~\cite{garingo2016}. These systems support remote rounds, specialist consultations, and family visits, cutting travel costs and reducing infection risks. However, studies also note several drawbacks. A review by Kristoffersson et al.~\cite{kristoffersson2013} found that telepresence robots, though functional, often feel impersonal and difficult to control. Clinicians report high cognitive load when navigating cluttered hospital spaces. Patients, especially the elderly and children, also find these robots less engaging than in-person visits due to limited physical embodiment and lack of expressive cues~\cite{sabanovic2013}. This “presence gap” limits their effectiveness beyond basic video conferencing, reducing their ability to offer comfort and reassurance during care.

\subsection{Relevant Literature Review}

\subsubsection{Overview of Existing Robots in Healthcare}
The hospital robotics market has matured significantly, with several manufacturers deploying autonomous mobile robots at scale. Table~\ref{tab:competitive} presents a comprehensive comparison of leading hospital robot platforms currently deployed in eastern markets.

% (Insert Table 2 here if needed)

\subsection{Success Stories: Autonomous Robots in Consumer and Industrial Sectors}

\subsubsection{Consumer Market: The Robotic Vacuum Revolution}
The most successful deployment of autonomous mobile robots to the general public has been robotic vacuum cleaners, which dominate the consumer robotics market without serious competition from traditional cleaning methods. The global robotic vacuum cleaner market has experienced explosive growth, with projections indicating expansion from \$9.07 billion in 2024 to \$31.79 billion by 2033. Global shipments reached 15.35 million units in the first half of recent years, demonstrating a 33\% year-over-year growth~\cite{scmp2024}.

The key to their success lies in the implementation of SLAM (Simultaneous Localization and Mapping) and localization technologies. Modern robotic vacuum cleaners employ LiDAR sensors, visual SLAM (VSLAM), and advanced mapping algorithms to create real-time maps of residential environments, enabling efficient path planning and obstacle avoidance~\cite{neato2025}. This technology has proven so reliable that robotic vacuums have achieved nearly 40\% of the service robot market share~\cite{trendforce2024}, demonstrating that autonomous navigation, when properly implemented, is mature enough for widespread consumer adoption.

\subsubsection{Industrial Market: Amazon’s Warehouse Robotics Dominance}
In the industrial sector, the most successful application of autonomous mobile robots has been in warehouse automation, with Amazon leading the field. Amazon has deployed over 750,000 to 1 million robots across its fulfillment network, making it the world’s largest fleet of industrial mobile robots~\cite{amazon2025}. These include various AMR systems such as:
\begin{itemize}
       \item \textbf{Proteus}: Amazon’s first fully autonomous mobile robot designed to work collaboratively with human workers in shared spaces, utilizing advanced sensors and safety features~\cite{robotsguide2023}.
       \item \textbf{Pegasus}: A drive unit optimized for transporting larger inventory pods, part of Amazon’s goods-to-person automation strategy~\cite{amazon2022}.
       \item \textbf{Hercules}: High-capacity drive units capable of lifting heavier loads within the warehouse environment.
       \item \textbf{Kiva Systems}: The foundational robotic platform (acquired in 2012) that initiated Amazon’s warehouse automation revolution, using QR code fiducials for localization~\cite{agv2024}.
\end{itemize}

\subsection{Core Technologies Breakdown}

\subsubsection{Navigation System}
The challenge of manual navigation highlighted in telepresence research is being addressed by autonomous mobile robots (AMRs). Systems like the Aethon TUG and the Swisslog Relay are now commonly used in hospitals to transport meals, lab samples, linens, and medications. These robots utilize technologies like LiDAR, SLAM (Simultaneous Localization and Mapping), and dynamic obstacle avoidance algorithms to navigate complex and crowded hospital corridors safely and efficiently.

The successful deployment of Autonomous Mobile Robots (AMRs) in operational environments presents significant challenges, particularly concerning safety and effective navigation within dynamic settings, which can lead to system failures and, notably, often results in user discomfort. A study examining three KMP AMRs successfully demonstrated the ability to navigate a static environment without collisions or safety events. However, this controlled environment exposed critical performance issues, including significant time variability between identical trips, poor docking accuracy required for recharging, and a relatively low average operating speed of 0.5 m/s substantially slower than the average human speed of 1.5 m/s. These robots, designed for predictable goods transport within product stores, contrast sharply with the objective of designing safe autonomous systems for highly dynamic and non-predictable environments.

Furthermore, the physical characteristics and cost of current industrial AMRs pose a barrier to wider adoption. The tested KMP robot models, with large dimensions (1.080m length, 0.700m height, and 0.640m width), necessitate complex components and extensive sensor arrays, contributing to a high unit cost. For instance, a basic navigation model like the KMP 600 is priced between \$80,000 and \$120,000, making it unaffordable for many potential users.

Accurate, real-time positioning is one of the key elements for successive AMR navigation, enabling both correct path-following and simultaneous map building. Real industrial robots often combine multiple solutions to enhance localization accuracy.

\paragraph{Radio Frequency (RF) Localization Technologies}
Various Radio Frequency (RF) technologies including Wi-Fi, RFID, Bluetooth Low Energy (BLE), ZigBee, and Ultra-Wideband (UWB) are employed, with selection criteria based on coverage, scalability, battery life, accuracy, and cost.

	extbf{Ultra-Wideband (UWB):} UWB solutions offer extended coverage and strong system scalability due to the presence of multiple nodes sharing the same spectrum. The battery life for a UWB sensor tag is excellent, comparable to other RF solutions, lasting for years. However, UWB is slightly more costly, with implementation over a 5000m$^2$ area estimated to cost up to \$20,000. Reported accuracy in real industrial scenarios varies widely, from 20 cm to more than 100 cm in studies on a 132m$^2$ area, and is subject to degradation in areas filled with obstacles or very large warehouses.

	extbf{Angle of Arrival (AoA):} Angle of Arrival (AoA), typically implemented using Bluetooth 5.1 or UWB Direction Finding, provides a low-energy localization solution that calculates a target's precise angle using a beacon (tag) and a robot-mounted antenna array, thereby eliminating the need for extensive static infrastructure. The reliance on Bluetooth Low Energy (BLE) grants AoA tags excellent battery life, lasting for years. While initial hardware costs are lower than other systems, the technology's main challenge is its inherent line-of-sight dependence and variable accuracy, which may necessitate multiple measurements. Advanced hybrid systems combining AoA with RSSI (Received Signal Strength Indication) have demonstrated high-precision potential in simulations, achieving mean accuracies of 12cm or less, even within complex, obstacle-filled environments.

\paragraph{SLAM-Based Sensor Fusion Techniques}
Beyond external RF systems, autonomous navigation heavily relies on Simultaneous Localization and Mapping (SLAM) approaches, often combining multiple sensors for robust performance.

	extbf{LiDAR, IMU, and Odometry Fusion:} A common hardware stack combines a 2D LiDAR (Light Detection and Ranging) with an IMU (Inertial Measurement Unit) and Odometry (e.g., wheel encoders). This sensor fusion is a powerful technique to achieve highly accurate real-time localization, as the IMU and Odometry provide constraints that correct for the motion distortion in LiDAR data and compensate for wheel slip, a major source of cumulative error in Odometry. This multi-sensor approach has been shown to significantly improve performance, achieving an improvement of nearly 27\% in rotation angle error and a reduction of over 67\% in position error compared to using a basic LiDAR SLAM system alone. The primary trade-off is the significant increase in system complexity and the computational resources required for processing the triple-stream data.

	extbf{Visual SLAM (VSLAM):} Visual SLAM utilizes cameras such as monocular, stereo, or RGB-D to achieve localization and map reconstruction. A key advantage of VSLAM is its significantly lower hardware cost and lighter weight compared to LiDAR-based methods, while also being capable of creating a richer environmental representation, including colors and textures. However, VSLAM is challenged by its reliance on environments with sufficient texture for reliable feature tracking and its difficulty in areas with repetitive patterns. Furthermore, early single-camera systems, such as ORB-SLAM 2.0, suffered from creating maps with unknown scales. The trend to mitigate these issues is the development of advanced multi-sensor frameworks, such as the hybrid Visual-Inertial-Ranging-LiDAR (VIRAL) SLAM system, which aims to enhance robustness in complex environments.

\subsubsection{Socially Assistive Robotics (SAR)}
In parallel, the field of Socially Assistive Robotics (SAR) has focused on the psycho-social aspects of patient care. These robots are designed not to perform physical tasks but to engage, encourage, and comfort patients through social interaction. For pediatric patients, robots like NAO and PARO (a therapeutic baby seal robot) have been shown to significantly reduce pain perception and anxiety during painful medical procedures such as vaccinations and chemotherapy~\cite{beran2013}. The robot acts as a coach and a distractor, creating a more positive and cooperative atmosphere. For elderly populations, SARs have been explored for companionship to combat loneliness, provide cognitive stimulation for dementia patients, and deliver medication reminders to improve treatment adherence~\cite{broekens2009}. These studies confirm that a robot’s ability to interact socially is a powerful tool for improving patient well-being and clinical outcomes. The primary limitation of SARs, however, is that they are typically specialized, non-mobile or semi-autonomous platforms that lack the high-fidelity telepresence capabilities needed for direct clinical consultation.

\subsection{Benefits of Autonomous Mobile Robots in Healthcare Settings}
Recent literature and deployment data reveal multiple evidence-based benefits of AMRs in hospital environments:
\begin{itemize}
       \item \textbf{Infection Control and Reduced Human Contact:} AMRs significantly minimize human-to-human contact, reducing cross-contamination risks—a critical advantage demonstrated during the COVID-19 pandemic~\cite{intel2024, frontiers2024}. Autonomous delivery robots eliminate the need for multiple staff members to handle potentially contaminated materials, reducing occupational exposure and conserving personal protective equipment. UV-C disinfection robots have been shown to achieve 99.9\% pathogen kill rates on environmental surfaces, helping reduce hospital-acquired infections (HAIs)~\cite{sciencedirect2021}.
       \item \textbf{Wayfinding and Visitor Guidance:} Intelligent hospital guidance robots provide accurate navigation assistance to visitors and patients, reducing confusion and staff burden~\cite{pmc2024nursing}. These robots can escort visitors to patient rooms, provide multilingual directions, and answer frequently asked questions about hospital facilities, allowing staff to focus on clinical care.
       \item \textbf{Material Transport and Delivery:} AMRs excel at routine delivery tasks including medications, meals, linens, laboratory specimens, and medical supplies~\cite{mdpi2024}. This automation frees nursing staff from non-clinical transport duties, allowing them to spend more time on direct patient care. Studies indicate that delivery robots can reduce staff time spent on logistics by up to 30\%.
       \item \textbf{Psychological Support and Patient Morale:} Research consistently demonstrates that pediatric patients exhibit positive emotional responses to robot interactions~\cite{beran2013}. Socially interactive robots reduce anxiety, provide distraction during painful procedures, and improve treatment compliance in children. Studies show that patients perceive robotic companions as non-judgmental, reducing psychological barriers to seeking help or expressing concerns~\cite{ahu2024}.
       \item \textbf{Autonomous Disinfection and Cleaning:} Cleaning and disinfection robots can be equipped with UV-C light systems at relatively low cost, providing continuous environmental sanitation~\cite{intel2024}. These robots operate during off-peak hours without human supervision, ensuring consistent hygiene standards while minimizing disruption to clinical workflows.
       \item \textbf{Telepresence and Virtual Presence:} AMRs with integrated telepresence capabilities enable remote consultations, family visits, and specialist rounds without requiring travel~\cite{garingo2016}. This technology proved invaluable during pandemic restrictions and continues to provide value for rural healthcare access and specialist consultations.
       \item \textbf{Waste Collection and Handling:} Some advanced AMR systems are deployed for biohazardous waste collection and disposal, reducing staff exposure to potentially infectious materials. In China and other Asian markets, robots have been successfully implemented for collecting and transporting medical waste from patient rooms to centralized disposal areas.
\end{itemize}

\subsection{Challenges in Hospital AMR Deployment}
Despite the significant benefits, several challenges remain:
\begin{itemize}
       \item \textbf{Cost and Return on Investment:} Hospital-grade AMRs represent substantial capital investments, typically ranging from \$50,000 to \$150,000 per unit depending on capabilities. Healthcare administrators require clear ROI demonstration through labor savings, efficiency gains, and improved outcomes to justify procurement~\cite{mdpi2024}.
       \item \textbf{System Integration Complexity:} Integrating AMRs with existing hospital information systems (HIS), electronic health records (EHR), elevator controls, automatic doors, and staff communication systems presents significant technical challenges. Interoperability standards are still evolving, requiring custom integration work for each deployment~\cite{frontiers2024}.
       \item \textbf{Safety, Privacy, and Trust:} Regulatory compliance (HIPAA, GDPR), patient data security, and physical safety in crowded environments remain critical concerns. Healthcare facilities must ensure that robots can safely navigate around patients, visitors, and medical equipment while maintaining patient privacy and data security~\cite{jmir2024}. Building trust among staff and patients requires careful change management and demonstration of reliability.
       \item \textbf{Regulatory and Liability Issues:} Medical device regulations, liability frameworks for autonomous systems, and insurance considerations create additional barriers to widespread adoption. Healthcare organizations must navigate complex approval processes and establish clear accountability frameworks for robot operations.
\end{itemize}

\subsection{Conclusion}
The literature presents a clear picture:
\begin{itemize}
       \item Telepresence robots are effective for remote diagnosis but are limited by manual control and an impersonal user experience.
       \item Socially assistive robots have proven their value in patient engagement and comfort but lack clinical telepresence and advanced mobility.
       \item Autonomous mobile robots have mastered navigation in hospitals but are used for logistics, not patient interaction.
       \item Conversational AI offers intelligent dialogue but lacks the physical embodiment needed for meaningful healthcare encounters.
       \item Consumer and industrial robotics have demonstrated that autonomous navigation technology is mature, reliable, and cost-effective for real-world deployment, as evidenced by millions of robotic vacuum cleaners in homes and over one million AMRs in Amazon warehouses.
       \item Current hospital AMRs are predominantly focused on material transport and disinfection, with limited integration of social interaction, advanced telepresence, and conversational AI capabilities.
\end{itemize}

The critical research gap lies at the intersection of these domains. There is a demonstrable need for a single, integrated platform that combines the clinical access of a telepresence robot, the patient-centric engagement of a SAR, the operational independence of an autonomous navigator, and the interactive intelligence of modern conversational AI. The proven success of autonomous navigation in consumer (robotic vacuums) and industrial (Amazon warehouses) applications demonstrates that the core technology is mature and ready for advanced healthcare applications. This project directly addresses this gap by designing a holistic system that synthesizes these strengths, creating a smart healthcare robot capable of not only connecting doctors and patients but also navigating its environment independently and engaging with patients in a helpful and empathetic manner. By learning from the successes of consumer robotics, warehouse automation, and existing hospital AMR deployments, this project aims to advance the state-of-the-art in healthcare robotics toward truly integrated, intelligent, and patient-centered care.

% --- END USER-REQUESTED LITERATURE REVIEW STRUCTURE ---
